\newcommand{\ModuleID}{EL-I.10}
\newcommand{\ModuleTitle}{Passive Voice and Deponent Verbs}
\newcommand{\ModuleStage}{Stage I — Foundations}
\newcommand{\ModuleUnit}{I.10}
\newcommand{\ModuleOutcome}{Distinguish true passives, deponents, participles, and forms of \Latin{fio}, while preserving subject, agency, and agreement}
\newcommand{\ModulePrerequisites}{EL-I.9 mastery: retrieve perfect stems and distinguish perfect-system time relations}
\newcommand{\ModuleExtensionPractice}{Worksheets I.38 and I.39 remain optional remediation or delayed-recall variants}
\newcommand{\ModuleNext}{EL-I.11, Pronouns and Prepositions}
\newcommand{\ModuleMastery}{Convert at least 18 of 20 active/passive forms without changing their other coordinates, build eight agreeing perfect passives, and classify at least nine of ten contextual passive-looking verbs as true passive, deponent, or \Latin{fio}.}
\newcommand{\ModuleDelayNote}{}
\newcommand{\ModuleLessonSource}{\CourseRoot/shared/content/volume1/all}
\newcommand{\ModuleExerciseSource}{\CourseRoot/shared/exercises/volume1/all}
\newcommand{\ModuleWorkedBank}{I}
\newcommand{\ModuleExerciseBank}{I}
\newcommand{\ModuleWorksheetVariants}{A,D}
\newcommand{\ModuleScopeNote}{This packet teaches the ordinary classification tree; dictionary identity controls whether a verb is deponent, and \Latin{fio} is not reduced to a regular passive of \Latin{facio}.}
\newcommand{\ModulePractice}{%
  \SubmoduleExtension
    {Present-system passive forms}
    {A true passive changes the relation between grammatical subject and action without erasing person, number, tense, or mood.}
    {Take three active forms from the lesson's paradigm and convert them to passive while holding every other coordinate constant. For each resulting clause, mark grammatical subject and any expressed personal agent or means.}
    {State which coordinates stay fixed in an active-to-passive form conversion.}
    {Person, number, tense, and mood remain fixed; voice changes. The former object may become grammatical subject in a clause transformation, while a personal agent is expressed with \Latin{a/ab} plus ablative and means ordinarily with the bare ablative.}
  \SubmoduleExtension
    {Perfect-system passives}
    {A perfect passive is a two-part agreement structure: participle plus finite \Latin{sum}.}
    {Use the lesson model to build a grid varying subject gender and number while holding the tense constant, then vary the tense while holding one subject constant. Draw one agreement arrow and one tense arrow per form.}
    {Identify which component carries gender and which carries person and tense.}
    {The perfect passive participle agrees with the subject in gender and number; the finite form of \Latin{sum} carries person, number, tense contribution, and mood. Both components are required for the complete finite group.}
  \SubmoduleExtension
    {Deponents}
    {A deponent's dictionary entry resolves the mismatch between passive morphology and active lexical meaning.}
    {From the lesson's deponent entries, make a card with present form, infinitive, perfect participial principal part, surface morphology, and active sense. Contrast one card with a true passive form of similar appearance.}
    {Explain why translating every \Latin{-tur} form as \English{is being} is unsafe.}
    {Each deponent card must preserve passive-form principal parts while assigning the active lexical sense shown in the lesson. The retrieval answer must require dictionary identity and context because \Latin{-tur} can occur in true passives and deponents.}
  \SubmoduleExtension
    {Handwritten study and controlled composition}
    {A four-way tree makes the evidence order explicit before translation.}
    {Classify the passive-looking forms in the lesson's final review as true finite passive, deponent finite verb, participle within a verbal group, or form of \Latin{fio}. For each, cite morphology, dictionary identity, and syntactic attachment.}
    {Recite the decision order without giving an English translation.}
    {A complete tree first identifies whether the form is finite or participial, then checks its dictionary entry, verbal companions, subject agreement, and clause role. \Latin{fio} retains its own lexical identity. Translation follows classification rather than deciding it.}
}
